% !TEX root = ../../ctfp-print.tex

\lettrine[lhang=0.17]{モ}{ナドについて説明したので}、双対性の恩恵を受けることができる。矢を逆にして反対の圏で作業するだけで、コモナドをただで得られる。

最も基本的なレベルでは、モナドはクライスリ矢の合成に関するものだったことを思い出してほしい：

\src{snippet01}
ここで \code{m} はモナドである関手だ。コモナドに文字 \code{w}（\code{m} を上下逆にしたもの）を使えば、コ・クライスリ矢を次の型の射として定義できる：

\src{snippet02}
コ・クライスリ矢に対する魚演算子の類似物は次のように定義される：

\src{snippet03}
コ・クライスリ矢が圏を形成するためには、恒等コ・クライスリ矢も必要であり、それは \code{extract} と呼ばれる：

\src{snippet04}
これは \code{return} の双対だ。結合律と左・右単位律も課す必要がある。まとめると、Haskell でコモナドを次のように定義できる：

\src{snippet05}
実際には、すぐに見るように、もう少し異なる基本要素を使う。

では、プログラミングにおいてコモナドは何の役に立つのだろうか？

\section{コモナドを使ったプログラミング}

モナドとコモナドを比較してみよう。モナドは \code{return} を使って値をコンテナに入れる手段を提供する。中に格納された値や値群にアクセスする手段は提供しない。もちろん、モナドを実装するデータ構造はその内容へのアクセスを提供するかもしれないが、それはボーナスと見なされる。モナドから値を取り出すための共通インターフェースは存在しない。そして、\code{IO} モナドが内容を決して公開しないことを誇りにしているという例を見てきた。

一方、コモナドは単一の値を取り出す手段を提供する。値を挿入する手段は提供しない。したがって、コモナドをコンテナとして考えたいなら、常に内容があらかじめ詰まっており、それを覗き見ることができる。

クライスリ矢が値を取ってある修飾された結果を生成する---文脈で修飾する---のと同様に、コ・クライスリ矢は値と全体の文脈を一緒に取って結果を生成する。それは\newterm{文脈計算}（contextual computation）の具現化だ。

\section{積コモナド}

リーダーモナドを覚えているだろうか？読み取り専用の環境 \code{e} へのアクセスが必要な計算を実装する問題に取り組むために導入した。そのような計算は次の形の純粋関数として表現できる：

\src{snippet06}
カリー化を使ってそれらをクライスリ矢に変換した：

\src{snippet07}
しかし、これらの関数はすでにコ・クライスリ矢の形をしていることに気づいてほしい。引数をより便利な関手の形に変形しよう：

\src{snippet08}
合成する矢に同じ環境を提供することで、合成演算子を簡単に定義できる：

\src{snippet09}
\code{extract} の実装は単に環境を無視する：

\src{snippet10}
当然ながら、積コモナドはリーダーモナドとまったく同じ計算を実行するために使える。ある意味で、環境のコモナド的な実装はより自然だ---「文脈における計算」の精神に従っている。一方、モナドは \code{do} 記法という便利な構文糖衣を持っている。

リーダーモナドと積コモナドの間の接続はより深く、リーダー関手が積関手の右随伴であるという事実に関係している。しかし一般に、コモナドはモナドとは異なる計算の概念をカバーする。後でさらに例を見ることになる。

\code{Product} コモナドをタプルやレコードを含む任意の積型に一般化することは簡単だ。

\section{合成の解剖}

双対化のプロセスを続けると、モナドのバインドと join を双対化することができる。あるいは、魚演算子の解剖を研究したモナドで使ったプロセスを繰り返すこともできる。このアプローチの方が示唆に富んでいるようだ。

出発点は、合成演算子は \code{w a} を取って \code{c} を生成するコ・クライスリ矢を生成しなければならないという認識だ。\code{c} を生成する唯一の方法は、型 \code{w b} の引数に二番目の関数を適用することだ：

\src{snippet11}
しかし、\code{g} に渡せる型 \code{w b} の値をどうやって生成できるか？手元にあるのは型 \code{w a} の引数と関数 \code{f :: w a -> b} だ。解決策は、バインドの双対を定義することで、それは extend と呼ばれる：

\src{snippet12}
\code{extend} を使って合成を実装できる：

\src{snippet13}
次に \code{extend} を解剖できるか？関数 \code{w a -> b} を引数 \code{w a} に適用すれば良いと思うかもしれないが、すぐに結果の \code{b} を \code{w b} に変換する方法がないことに気づく。コモナドは値を持ち上げる手段を提供しないことを覚えておいてほしい。この時点で、モナドの類似した構成では \code{fmap} を使った。ここで \code{fmap} を使う唯一の方法は、型 \code{w (w a)} のものが手元にある場合だ。\code{w a} を\\ \code{w (w a)} に変換できれば良いのだが。そして便利なことに、それはまさに \code{join} の双対になる。それを \code{duplicate} と呼ぶ：

\src{snippet14}
したがって、モナドの定義と同様に、コモナドには三つの等価な定義がある：コ・クライスリ矢を使うもの、\code{extend} を使うもの、\code{duplicate} を使うものだ。\\ \code{Control.Comonad} ライブラリから直接取った Haskell の定義がこちらだ：

\src{snippet15}
\code{duplicate} による \code{extend} のデフォルト実装とその逆が提供されているので、どちらか一方をオーバーライドするだけで良い。

これらの関数の直観は、一般にコモナドは型 \code{a} の値が詰まったコンテナ（積コモナドは値が一つだけの特殊なケースだった）として考えられるというアイデアに基づいている。「現在の」値という概念があり、それは \code{extract} を通じて簡単にアクセスできる。コ・クライスリ矢は現在の値に焦点を当てた何らかの計算を実行するが、周囲の値すべてにアクセスできる。コンウェイのライフゲームを考えてみよう。各セルは値（通常は \code{True} または \code{False} だけ）を含む。ライフゲームに対応するコモナドは「現在の」セルに焦点を当てたセルのグリッドになる。

では \code{duplicate} は何をするのか？コモナドのコンテナ \code{w a} を取り、コンテナのコンテナ \code{w (w a)} を生成する。アイデアは、これらのコンテナそれぞれが \code{w a} 内の異なる \code{a} に焦点を当てているということだ。ライフゲームでは、グリッドのグリッドを得ることになり、外側のグリッドの各セルは異なるセルに焦点を当てた内側のグリッドを含む。

次に \code{extend} を見てみよう。コ・クライスリ矢とコモナドのコンテナ \code{w a}（\code{a} が詰まっている）を取る。これらの \code{a} すべてに計算を適用し、\code{b} に置き換える。結果は \code{b} が詰まったコモナドのコンテナだ。\code{extend} は一つの \code{a} から別の \code{a} へと焦点をシフトさせ、それぞれにコ・クライスリ矢を順番に適用することでこれを行う。ライフゲームでは、コ・クライスリ矢は現在のセルの新しい状態を計算する。そのために文脈---おそらく最近傍のセル---を参照する。\code{extend} のデフォルト実装がこのプロセスを説明している。まず \code{duplicate} を呼んで可能なすべての焦点を生成し、次にそれぞれに \code{f} を適用する。

\section{ストリームコモナド}

コンテナの一つの要素から別の要素へと焦点をシフトするこのプロセスは、無限ストリームの例で最もよく説明される。そのようなストリームはリストとほぼ同じだが、空のコンストラクタを持たない：

\src{snippet16}
これは自明に \code{Functor} だ：

\src{snippet17}
ストリームの焦点はその最初の要素なので、\code{extract} の実装はこうなる：

\src{snippet18}
\code{duplicate} はストリームのストリームを生成し、それぞれが異なる要素に焦点を当てる。

\src{snippet19}
最初の要素は元のストリームで、二番目の要素は元のストリームの尾部で、三番目の要素はその尾部、以下同様に無限に続く。

完全なインスタンスはこうなる：

\src{snippet20}
これはストリームを見る非常に関数的な方法だ。命令型言語では、おそらく一つの位置ずつストリームを進める \code{advance} メソッドから始めるだろう。ここでは、\code{duplicate} がすべてのシフトされたストリームを一度に生成する。Haskell の遅延評価がこれを可能にし、望ましくさえある。もちろん、\code{Stream} を実用的にするには \code{advance} の類似物も実装する：

\src{snippet21}
しかし、それはコモナドのインターフェースの一部ではない。

デジタル信号処理の経験があれば、ストリームに対するコ・クライスリ矢はデジタルフィルターに過ぎず、\code{extend} がフィルタリングされたストリームを生成することがすぐにわかるだろう。

簡単な例として、移動平均フィルターを実装してみよう。ストリームの \code{n} 要素を合計する関数はこうなる：

\src{snippet22}
ストリームの最初の \code{n} 要素の平均を計算する関数はこうなる：

\src{snippet23}
部分適用された \code{average n} はコ・クライスリ矢なので、ストリーム全体に \code{extend} できる：

\src{snippet24}
結果は移動平均のストリームだ。

ストリームは一方向の一次元コモナドの例だ。双方向にしたり、二次元以上に拡張したりすることも簡単だ。

\section{圏論的コモナド}

圏論でコモナドを定義することは、双対性における単純な演習だ。モナドと同様に、自己関手 $T$ から始める。モナドを定義する二つの自然変換 $\eta$ と $\mu$ は、コモナドでは単純に逆になる：
\begin{align*}
  \varepsilon & \Colon T \to I   \\
  \delta      & \Colon T \to T^2
\end{align*}
これらの変換の成分は \code{extract} と \code{duplicate} に対応する。コモナド則はモナド則の鏡像だ。大きな驚きはない。

そして、随伴からモナドを導く方法がある。双対性は随伴を逆にする：左随伴は右随伴になり、逆も然りだ。そして、合成 $R \circ L$ がモナドを定義するので、$L \circ R$ はコモナドを定義しなければならない。随伴の余単位：
\[\varepsilon \Colon L \circ R \to I\]
は確かにコモナドの定義に見られる $\varepsilon$ と同じだ---あるいは成分では、Haskell の \code{extract} だ。随伴の単位も使える：
\[\eta \Colon I \to R \circ L\]
これを使って $L \circ R$ の途中に $R \circ L$ を挿入し、$L \circ R \circ L \circ R$ を生成できる。$T$ から $T^2$ を作ることで $\delta$ が定義され、それでコモナドの定義が完成する。

また、モナドはモノイドであることも見てきた。この陳述の双対はコモノイドの使用を必要とするが、コモノイドとは何だろうか？単一対象圏としてのモノイドの元の定義は何も面白いものに双対化されない。すべての自己射の方向を逆にすると、別のモノイドが得られる。しかし、モナドへのアプローチでは、モノイダル圏における対象としてのモノイドのより一般的な定義を使った。その構成は二つの射に基づいていた：
\begin{align*}
  \mu  & \Colon m \otimes m \to m \\
  \eta & \Colon i \to m
\end{align*}
これらの射の逆転によって、モノイダル圏における\newterm{コモノイド}（comonoid）が生まれる：
\begin{align*}
  \delta      & \Colon m \to m \otimes m \\
  \varepsilon & \Colon m \to i
\end{align*}
Haskell でコモノイドの定義を書くことができる：

\src{snippet25}
しかし、それはかなり自明だ。明らかに \code{destroy} はその引数を無視する。

\src{snippet26}
\code{split} は単に一対の関数だ：

\src{snippet27}
ここで、モノイドの単位律の双対であるコモノイド則を考えよう。

\src{snippet28}
ここで、\code{lambda} と \code{rho} はそれぞれ左単位子と右単位子だ（\hyperref[monads-categorically]{モノイダル圏}の定義を参照）。定義を代入すると：

\src{snippet29}
これにより \code{g = id} が証明される。同様に、二番目の法則は \code{f = id} に展開される。結論として：

\src{snippet30}
これは Haskell（そして一般に $\Set$ 圏）では、すべての対象が自明なコモノイドであることを示している。

幸いにも、コモノイドを定義するためのより興味深いモノイダル圏が他にある。その一つが自己関手の圏だ。そして、モナドが自己関手の圏におけるモノイドであるのと同様に、

\begin{quote}
  コモナドは自己関手の圏におけるコモノイドだ。
\end{quote}

\section{ストアコモナド}

コモナドのもう一つの重要な例は状態モナドの双対だ。それはコ状態コモナド、あるいは別名\newterm{ストアコモナド}（store comonad）と呼ばれる。

状態モナドは指数を定義する随伴によって生成されることを前に見た：
\begin{align*}
  L z & = z\times{}s      \\
  R a & = s \Rightarrow a
\end{align*}
同じ随伴を使ってコ状態コモナドを定義する。コモナドは合成 $L \circ R$ によって定義される：
\[L (R a) = (s \Rightarrow a)\times{}s\]
これを Haskell に翻訳するにあたって、左の \code{Product} 関手と右の \code{Reader} 関手の間の随伴から始める。\code{Reader} の後に \code{Product} を合成することは次の定義と等価だ：

\src{snippet31}
対象 $a$ における随伴の余単位は射：
\[\varepsilon_a \Colon ((s \Rightarrow a)\times{}s) \to a\]
あるいは Haskell 記法では：

\src{snippet32}
これが \code{extract} になる：

\src{snippet33}
随伴の単位：

\src{snippet34}
は部分適用されたデータコンストラクタとして書き直せる：

\src{snippet35}
$\delta$、すなわち \code{duplicate} を水平合成として構成する：
\begin{align*}
  \delta & \Colon L \circ R \to L \circ R \circ L \circ R \\
  \delta & = L \circ \eta \circ R
\end{align*}
$\eta$、すなわち \code{Store f} を最も左の $L$（\code{Product} 関手）を通じて忍び込ませる必要がある。これはペアの左の成分に $\eta$ で作用することを意味する（\code{Product} に対する \code{fmap} がすることだ）。こうなる：

\src{snippet36}
（$\delta$ の公式では、$L$ と $R$ は成分が恒等射である恒等自然変換を表すことを覚えておいてほしい。）

\code{Store} コモナドの完全な定義がこちらだ：

\src{snippet37}
\code{Store} の \code{Reader} 部分を型 \code{s} の要素を使ってキーが付けられた \code{a} の一般化コンテナとして考えることができる。例えば、\code{s} が \code{Int} の場合、\code{Reader Int a} は \code{a} の無限双方向ストリームだ。\code{Store} はこのコンテナをキー型の値と対にする。例えば、\code{Reader Int a} は \code{Int} と対にされる。この場合、\code{extract} はこの整数を使って無限ストリームにインデックスを付ける。\code{Store} の二番目の成分を現在の位置と考えることができる。

この例を続けると、\code{duplicate} は \code{Int} でインデックスが付けられた新しい無限ストリームを作成する。このストリームはその要素としてストリームを含む。特に、現在の位置では元のストリームを含む。しかし、別の \code{Int}（正または負）をキーとして使うと、その新しいインデックスに位置するシフトされたストリームを得る。

一般に、\code{extract} が \code{duplicate} された \code{Store} に作用すると元の \code{Store} が生成されることを自分で確認できる（実際、コモナドの恒等律は \code{extract . duplicate = id} と述べている）。

\code{Store} コモナドは \code{lens} ライブラリの理論的基礎として重要な役割を果たす。概念的に、\code{Store s a} コモナドは型 \code{s} をインデックスとして使ってデータ型 \code{a} の特定の部分構造に「焦点を当てる」（レンズのような）というアイデアをカプセル化する。特に、型の関数：

\src{snippet38}
は一対の関数と等価だ：

\src{snippet39}
\code{a} が積型の場合、\code{set} は \code{a} の内側の型 \code{s} のフィールドを設定し、\code{a} の変更されたバージョンを返すものとして実装できる。同様に、\code{get} は \code{a} から \code{s} フィールドの値を読み取るものとして実装できる。次のセクションでこれらのアイデアをさらに探求する。

\section{練習問題}

\begin{enumerate}
  \tightlist
  \item
        \code{Store} コモナドを使ってコンウェイのライフゲームを実装せよ。
        ヒント：\code{s} にどの型を選ぶか？
\end{enumerate}
